% ============================================================================
%  pipeline_flow.tex -- how the Adaptive ML Experiment Agent answers a
%  research question, explained for a reader who has never seen the code.
%
%  Part 1: the mental model.  Part 2: the setup both modes share.
%  Part 3: Selection Mode end to end.  Part 4: Effect Mode end to end.
%  Part 5: the two modes side by side.
%
%  Every number marked "real run" comes from running the current code on
%  diabetes_risk.csv (real ingestion, training and statistics; the LLM
%  replies were scripted, as in the test suite). Numbers marked
%  "illustration" only show the mechanism.
%
%  Compile:  pdflatex pipeline_flow.tex    or    tectonic pipeline_flow.tex
% ============================================================================
\documentclass[11pt]{article}
\usepackage[a4paper,margin=1.6cm]{geometry}
\usepackage[T1]{fontenc}
\usepackage[scaled=0.94]{helvet}
\renewcommand{\familydefault}{\sfdefault}
\usepackage[scaled=0.96]{inconsolata}
\usepackage{amsmath,amssymb}
\usepackage{xcolor}
\usepackage{tikz}
\usetikzlibrary{arrows.meta,positioning,calc,fit,backgrounds}
\usepackage[most]{tcolorbox}
\usepackage{booktabs,tabularx,array}
\usepackage{enumitem}
\usepackage{microtype}
\usepackage{ragged2e}
\usepackage{needspace}

\pagestyle{empty}
\setlength{\parindent}{0pt}
\setlength{\parskip}{0pt}

% ---- palette ---------------------------------------------------------------
% ink/rule/soft: text and structure. One colour per kind of actor:
%   code (grey), Planner LLM (blue), Recommender LLM (violet),
%   validation (green), sealed test (red).
\definecolor{ink}{HTML}{1B1F24}
\definecolor{rule}{HTML}{BCC3CC}
\definecolor{soft}{HTML}{F2F4F6}
\definecolor{plan}{HTML}{1F4E79}  \definecolor{planbg}{HTML}{DCE8F4}
\definecolor{rec}{HTML}{6A3D9A}   \definecolor{recbg}{HTML}{ECE4F5}
\definecolor{val}{HTML}{2E7D4F}   \definecolor{valbg}{HTML}{E2F1E7}
\definecolor{tst}{HTML}{B03A2E}   \definecolor{tstbg}{HTML}{F8E4E1}
\color{ink}

\setlist[enumerate]{nosep,leftmargin=1.4em,label=\textbf{\arabic*.},itemsep=2pt}
\setlist[itemize]{nosep,leftmargin=1.1em,label=--,itemsep=2pt}
\newcolumntype{Y}{>{\RaggedRight\arraybackslash}X}
\renewcommand{\arraystretch}{1.2}

% ---- building blocks ------------------------------------------------------
\newcommand{\dat}[1]{\texttt{#1}}
\newcommand{\gap}{\par\vspace{7pt}}
\newcommand{\lbl}[1]{{\footnotesize\bfseries\color{plan}\MakeUppercase{#1}}\par\vspace{3pt}}
\newcommand{\partx}[2]{\par\needspace{8\baselineskip}\vspace{6pt}%
  \noindent{\large\bfseries #1\enspace #2}\par\nobreak\vspace{3pt}\nobreak
  {\color{rule}\hrule height 0.8pt}\nobreak\vspace{8pt}}
\newcommand{\step}[3]{\par\needspace{5\baselineskip}\vspace{9pt}%
  \noindent{\bfseries\color{plan}#1}\enspace{\bfseries #2}\hfill #3\par\nobreak\vspace{4pt}}
\newcommand{\note}[1]{{\small\color{black!62}#1}}
\newcommand{\tagcode}{\colorbox{soft}{\footnotesize\strut deterministic code}}
\newcommand{\tagplan}{\colorbox{planbg}{\footnotesize\strut\color{plan}\textbf{Planner LLM}}}
\newcommand{\tagrec}{\colorbox{recbg}{\footnotesize\strut\color{rec}\textbf{Recommender LLM}}}
\newcommand{\tagval}{\colorbox{valbg}{\footnotesize\strut\color{val}\textbf{validation}}}
\newcommand{\tagtst}{\colorbox{tstbg}{\footnotesize\strut\color{tst}\textbf{sealed test}}}

\newtcolorbox{panel}[1][]{%
  enhanced, colback=soft, colframe=soft, boxrule=0pt, arc=1.5pt,
  left=6pt, right=6pt, top=4pt, bottom=4pt, before skip=3pt, after skip=4pt,
  before upper=\RaggedRight, #1}
\newtcolorbox{keybox}[1][]{%
  enhanced, colback=white, colframe=plan, boxrule=0.7pt, arc=1.5pt,
  left=7pt, right=7pt, top=5pt, bottom=5pt, before skip=4pt, after skip=6pt,
  before upper=\RaggedRight, #1}

% ---- TikZ styles shared by every diagram ----------------------------------
\tikzset{
  box/.style={draw=rule, fill=white, rounded corners=2pt, line width=0.6pt,
    align=center, font=\small, inner sep=4pt, minimum height=7mm},
  code/.style={box},
  planner/.style={box, draw=plan, fill=planbg, line width=1pt},
  recom/.style={box, draw=rec, fill=recbg, line width=1pt},
  valid/.style={box, draw=val, fill=valbg, line width=1pt},
  test/.style={box, draw=tst, fill=tstbg, line width=1pt},
  io/.style={box, fill=soft},
  q/.style={box, fill=soft, font=\small\itshape},
  arr/.style={-{Stealth[length=2mm]}, line width=0.7pt, black!55},
  lab/.style={font=\footnotesize\color{black!60}},
  run/.style={draw=rule, fill=white, rounded corners=1.5pt, font=\scriptsize\ttfamily,
    inner sep=2pt, minimum height=4.5mm},
}

\begin{document}
\RaggedRight

% ============================================================================
%  PART 1 -- THE MENTAL MODEL
% ============================================================================
{\fontsize{18}{22}\selectfont\bfseries How the Adaptive ML Experiment Agent answers a question}\par
\vspace{3pt}
{\color{black!60}A research question and a CSV go in; a statistically honest answer comes out.
This page is the whole idea. The rest of the document opens each box.}\par
\vspace{8pt}

\begin{center}
\begin{tikzpicture}[node distance=5mm]
  % main column
  \node[io, minimum width=6.6cm] (q) {\textbf{Question + dataset}};
  \node[code, minimum width=6.6cm, below=of q] (ing) {Ingest \& profile the dataset};
  \node[planner, minimum width=6.6cm, below=of ing] (pl) {\textbf{Planner LLM}: choose the initial experiment design};
  \node[code, minimum width=6.6cm, below=of pl] (des) {Build candidates, then seeds $\rightarrow$ runs};
  \node[code, minimum width=6.6cm, below=of des] (run) {Run experiments (train every run)};
  \node[valid, minimum width=6.6cm, below=of run] (va) {\textbf{Validation analysis}: averages, paired bootstrap};
  \node[code, minimum width=6.6cm, below=of va, minimum height=9mm] (dec) {\textbf{Decide}\\[-1pt]{\footnotesize code checks the stop rules first}};
  \node[code, minimum width=6.6cm, below=10mm of dec] (fin) {Finalize: choose the final comparisons};
  \node[test, minimum width=6.6cm, below=of fin] (te) {\textbf{Test set opened once}: confirmatory comparisons};
  \node[code, minimum width=6.6cm, below=of te] (rep) {Report \;{\footnotesize(+ LLM explanation of the numbers)}};
  \foreach \a/\b in {q/ing, ing/pl, pl/des, des/run, run/va, va/dec, fin/te, te/rep}
    \draw[arr] (\a) -- (\b);
  \draw[arr] (dec) -- node[right, lab] {conclude} (fin);
  % loop column (right)
  \node[recom, text width=3.4cm, right=12mm of dec] (rc) {\textbf{Recommender LLM}\\[-1pt]{\footnotesize proposes one refinement}};
  \node[code, text width=3.4cm, right=12mm of run] (vb) {Code validates it and builds new candidates};
  \draw[arr] (dec) -- node[above, lab] {refine} (rc);
  \draw[arr] (rc) -- (vb);
  \draw[arr] (vb) -- node[above, lab] {new runs} (run);
  \draw[arr] (rc.south) |- node[pos=0.25, right, lab] {conclude} (fin.east);
  % sealed test (left)
  \node[test, text width=2.7cm, left=12mm of va, font=\footnotesize] (vault) {\textbf{TEST rows}\\scored at training time,\\ \textbf{sealed}};
  \draw[arr, tst, dashed] (run.west) -| node[pos=0.25, above, lab] {stored} (vault.north);
  \draw[arr, tst, dashed] (vault.south) |- node[pos=0.75, above, lab] {read once} (te.west);
  % the loop label
  \node[lab, right=1mm of vb.north east, anchor=south east, yshift=1mm] {the adaptive loop};
\end{tikzpicture}
\end{center}

\vspace{4pt}
\begin{tabularx}{\linewidth}{@{}l Y@{}}
\tagcode & Every rule and every number: candidates, seeds, training, statistics, stop rules, validation of LLM proposals.\\
\tagplan & One call. Reads the question and dataset \emph{counts}; returns the initial design. Trains nothing, computes nothing, never sees a data row.\\
\tagrec & At most one call per round. Reads a compact summary of validation results; \emph{recommends} one refinement or stopping. Code decides whether it is allowed.\\
\tagval & The only data the adaptive loop ever looks at. Paired bootstrap happens here, every round.\\
\tagtst & Scored when each model is trained, then sealed. Read exactly once, after the loop has stopped and the comparisons are fixed.\\
\end{tabularx}

\vspace{6pt}
\begin{keybox}
\textbf{Two kinds of question, one skeleton.}
\textbf{Selection Mode} asks \emph{which model family is best?} --- the initial design is several model
families, and refinement tunes one knob of a family still in the race.
\textbf{Effect Mode} asks \emph{does factor X matter?} --- the initial design is several levels of one
factor, and refinement only adds more levels of that same factor. Training, seeds, the bootstrap and the
sealed test set work identically in both.
\end{keybox}

% ============================================================================
%  PART 2 -- COMMON SETUP
% ============================================================================
\newpage
\partx{2}{Common setup --- the same for both modes}

\step{A}{Ingest \& profile}\tagcode

\begin{minipage}[t]{0.52\linewidth}
\vspace{0pt}
\begin{itemize}
\item \textbf{In:} the CSV, the target column the user names (never guessed), the question.
\item Rows with any missing value are dropped (no imputation).
\item Task type: text target, or $\le 20$ whole-number values $\rightarrow$ classification;
      otherwise regression.
\item Categorical columns are one-hot encoded.
\item One fixed split for every run, ever: \textbf{70\,\% train / 15\,\% validation / 15\,\% test}.
\item \textbf{Out:} a \dat{DatasetProfile} (counts, classes, split seed).
\end{itemize}
\end{minipage}\hfill
\begin{minipage}[t]{0.44\linewidth}
\vspace{0pt}
\begin{panel}
\small\textbf{Real run --- \dat{diabetes\_risk.csv}}\par\vspace{2pt}
15\,000 rows $\rightarrow$ 10\,500 after dropping missing values\\
target \dat{diabetes\_risk}: Low / Moderate / High $\rightarrow$ classification\\
18 feature columns $\rightarrow$ 49 after one-hot\\
split: 7\,350 train / 1\,575 validation / 1\,575 test\\
always predicting ``Low'' scores 60.0\,\% (the sanity floor)
\end{panel}
\end{minipage}

\step{B}{Planner LLM --- the initial design}\tagplan

\begin{center}
\begin{tikzpicture}[node distance=4mm]
  \node[io, text width=4.6cm, font=\footnotesize] (in) {\textbf{Receives}\\question text\\+ dataset \emph{counts}: task, target name,\\rows, column counts, class sizes};
  \node[planner, right=8mm of in, text width=3.1cm] (pl) {\textbf{Planner LLM}\\{\footnotesize one call, strict JSON}};
  \node[code, right=8mm of pl, text width=3.9cm, font=\footnotesize] (ck) {\textbf{Code checks the design}\\families / factor exist,\\levels distinct and in range\\(a violation is sent back once)};
  \draw[arr] (in) -- (pl); \draw[arr] (pl) -- (ck);
  \coordinate (j) at ($(pl.south)+(0,-9mm)$);
  \node[box, draw=plan, line width=1pt, text width=5.4cm, font=\footnotesize, anchor=north] (sel)
    at ($(j)+(-3.3,-5mm)$)
    {\textbf{IF mode = SELECTION}\\2--4 model families to compare\\$\rightarrow$ Part 3};
  \node[box, draw=plan, line width=1pt, text width=5.4cm, font=\footnotesize, anchor=north] (eff)
    at ($(j)+(3.3,-5mm)$)
    {\textbf{IF mode = EFFECT}\\one factor, 2--5 levels, a reference level,\\the base configuration $\rightarrow$ Part 4};
  \draw[line width=0.7pt, black!55] (ck.south) |- (j);
  \draw[arr] (j) -| (sel.north);
  \draw[arr] (j) -| (eff.north);
\end{tikzpicture}
\end{center}

\begin{minipage}[t]{0.48\linewidth}
\textbf{The Planner does NOT}
\begin{itemize}
\item train models or run experiments,
\item calculate any statistic,
\item see individual data rows or column names.
\end{itemize}
\end{minipage}\hfill
\begin{minipage}[t]{0.48\linewidth}
\textbf{The four model families it can choose from}\\[2pt]
{\small logistic / linear regression (\dat{linear\_baseline}), decision tree, random forest, MLP
(one hidden layer). Each has fixed defaults and allowed ranges in a code registry.}
\end{minipage}

\vspace{8pt}
\begin{keybox}
\textbf{Candidate $\neq$ run.} A \textbf{candidate} is one configuration --- model family + hyperparameters
+ preprocessing --- \emph{without} a random seed. A \textbf{run} is one actual training of a candidate
with one seed. Random forests and MLPs depend on the seed, so each of their candidates is trained
\textbf{3 times} (seeds 0, 1, 2). Logistic regression and the decision tree do not, so they are trained
\textbf{once}. Every statistic is computed per \emph{candidate}, after averaging its runs.
\end{keybox}

% ============================================================================
%  PART 3 -- SELECTION MODE
% ============================================================================
\newpage
\partx{3}{Selection Mode --- ``which model family is best?''}

{\small Running example (real run): \emph{``Which model performs best for classification on this
dataset?''}; the Planner chooses all four families.}

\begin{center}
\begin{tikzpicture}[node distance=4.2mm]
  \node[code, minimum width=6.4cm] (s1) {S1\; one root candidate per family};
  \node[code, minimum width=6.4cm, below=of s1] (s2) {S2\; seeds $\rightarrow$ runs};
  \node[code, minimum width=6.4cm, below=of s2] (s3) {S3\; train; score every val \& test row};
  \node[valid, minimum width=6.4cm, below=of s3] (s4) {S4\; average seeds, find the leader};
  \node[valid, minimum width=6.4cm, below=of s4] (s5) {S5\; paired bootstrap $\rightarrow$ contenders};
  \node[code, minimum width=6.4cm, below=of s5] (s6) {S6\; stop rules (code)};
  \node[code, minimum width=6.4cm, below=9mm of s6] (s8) {S8\; choose winner $W$ + runner-up $R$};
  \node[test, minimum width=6.4cm, below=of s8] (s9) {S8\; open test: $W$ vs reference, $W$ vs $R$};
  \foreach \a/\b in {s1/s2,s2/s3,s3/s4,s4/s5,s5/s6,s8/s9} \draw[arr] (\a) -- (\b);
  \draw[arr] (s6) -- node[right, lab] {stop} (s8);
  \node[recom, text width=3.6cm, right=12mm of s6] (r) {S6\; \textbf{Recommender}:\\{\footnotesize one contender, one knob,\\1--3 values}};
  \node[code, text width=3.6cm, right=12mm of s2] (b) {S7\; code validates,\\builds 1--3 children};
  \draw[arr] (s6) -- node[above, lab] {continue} (r);
  \draw[arr] (r) -- (b);
  \draw[arr] (b) -- node[above, lab] {new} (s2);
  \draw[arr] (r.south) |- node[pos=0.25, right, lab] {conclude} (s8.east);
\end{tikzpicture}
\end{center}

% ---------------------------------------------------------------- S1 + S2
\step{S1}{Create root candidates}\tagcode

Each chosen family gets exactly \textbf{one} root candidate at its registry defaults. The simplest chosen
family (order: linear $<$ tree $<$ forest $<$ MLP) becomes the \textbf{reference} --- the baseline every
other candidate is measured against.

\begin{panel}
\small\ttfamily
c1 = linear\_baseline \hfill \textsf{reference}\\
c2 = decision\_tree \ (max\_depth 8)\\
c3 = random\_forest \ (200 trees)\\
c4 = mlp \ (hidden 64, 20 epochs, features normalized)
\end{panel}

\step{S2}{Expand candidates into runs}\tagcode

\begin{minipage}[t]{0.47\linewidth}
\vspace{0pt}
\begin{tikzpicture}[x=1cm, y=-0.62cm]
  \node[code, anchor=west, font=\footnotesize, minimum height=5mm] (a) at (0,0) {c1 linear};
  \node[code, anchor=west, font=\footnotesize, minimum height=5mm] (b) at (0,1) {c2 tree};
  \node[code, anchor=west, font=\footnotesize, minimum height=5mm] (c) at (0,2.6) {c3 forest};
  \node[code, anchor=west, font=\footnotesize, minimum height=5mm] (d) at (0,5.6) {c4 MLP};
  \node[run, anchor=west] (a0) at (3.0,0) {run: seed 0};
  \node[run, anchor=west] (b0) at (3.0,1) {run: seed 0};
  \foreach \s/\y in {0/1.8,1/2.6,2/3.4} \node[run, anchor=west] (c\s) at (3.0,\y) {run: seed \s};
  \foreach \s/\y in {0/4.8,1/5.6,2/6.4} \node[run, anchor=west] (d\s) at (3.0,\y) {run: seed \s};
  \draw[arr] (a) -- (a0); \draw[arr] (b) -- (b0);
  \foreach \s in {0,1,2} { \draw[arr] (c.east) -- (c\s.west); \draw[arr] (d.east) -- (d\s.west); }
  \node[lab, anchor=west] at (4.75,0) {seed-independent};
  \node[lab, anchor=west] at (4.75,1) {seed-independent};
\end{tikzpicture}
\end{minipage}\hfill
\begin{minipage}[t]{0.49\linewidth}
\vspace{2pt}
\begin{panel}
\textbf{4 candidates $\rightarrow$ 8 actual training runs.}\par\vspace{3pt}
\small A seed changes only randomness \emph{inside} training (weight initialisation, dropout masks,
batch order, a forest's bootstrap samples). It never changes which rows are train, validation or test.
\end{panel}
\end{minipage}

% ---------------------------------------------------------------- S3
\needspace{17\baselineskip}
\step{S3}{Train --- and score every row}\tagcode

\begin{center}
\begin{tikzpicture}
  \node[code, minimum width=7.6cm, minimum height=8mm, anchor=west] (tr) at (0,0) {\textbf{TRAIN} 70\,\% (7\,350 rows) $\rightarrow$ fit the model};
  \node[valid, minimum width=3.5cm, minimum height=8mm, anchor=west, right=1mm of tr] (vl) {\textbf{VALIDATION} 15\,\%};
  \node[test, minimum width=3.5cm, minimum height=8mm, anchor=west, right=1mm of vl] (ts) {\textbf{TEST} 15\,\%};
  \node[font=\footnotesize\ttfamily, below=2mm of vl] (vs) {1 0 1 1 \dots{} \textsf{1\,575 rows}};
  \node[font=\footnotesize\ttfamily, below=2mm of ts] (tss) {0 1 1 1 \dots{} \textsf{1\,575 rows}};
  \draw[arr, val] (vl) -- (vs); \draw[arr, tst] (ts) -- (tss);
  \node[font=\footnotesize\color{val}\bfseries, below=0.5mm of vs] {available to the analysis};
  \node[font=\footnotesize\color{tst}\bfseries, below=0.5mm of tss] {stored, \textbf{sealed}};
\end{tikzpicture}
\end{center}

\begin{minipage}[t]{0.5\linewidth}
Each run scores \textbf{every} validation row and \textbf{every} test row --- not just one accuracy number:
\[
\text{score}(i) =
\begin{cases}
1 \text{ correct},\; 0 \text{ wrong} & \text{classification}\\
(\hat y_i - y_i)^2 & \text{regression}
\end{cases}
\]
\end{minipage}\hfill
\begin{minipage}[t]{0.46\linewidth}
\vspace{2pt}
{\small It also records the same metric on the \emph{training} rows (to spot overfitting), the training
loss and the time. A run that crashes or produces NaN/Inf is marked \dat{failed}; every other run counts
as evidence. Each run is saved as it finishes.}
\end{minipage}

% ---------------------------------------------------------------- S4
\step{S4}{Validation analysis --- average the seeds, find the leader}\tagval

\begin{minipage}[t]{0.5\linewidth}
\vspace{0pt}
\textbf{First, seeds are averaged row by row} \note{(illustration)}\par\vspace{3pt}
{\small
\begin{tabular}{@{}l c c c c c@{}}
 & row 1 & row 2 & row 3 & row 4 & \dots\\ \midrule
forest seed 0 & 1 & 0 & 1 & 1 & \\
forest seed 1 & 1 & 1 & 1 & 0 & \\
forest seed 2 & 0 & 1 & 1 & 1 & \\ \midrule
\textbf{c3 forest} & \textbf{.67} & \textbf{.67} & \textbf{1} & \textbf{.67} & \\
\end{tabular}}
\par\vspace{4pt}
{\small $\Rightarrow$ one validation score \emph{per row} per candidate. Its mean over the 1\,575 rows is
the candidate's validation accuracy.}
\end{minipage}\hfill
\begin{minipage}[t]{0.45\linewidth}
\vspace{0pt}
\textbf{Then each candidate gets one number} \note{(real run)}\par\vspace{3pt}
{\small
\begin{tabular}{@{}l r r@{}}
\toprule
candidate & validation & train\\ \midrule
c1 linear (reference) & \textbf{79.05\,\%} & 80.0\,\%\\
c2 tree & 75.75\,\% & 83.4\,\%\\
c3 forest & 78.20\,\% & 100.0\,\%\\
c4 MLP & 77.40\,\% & 85.0\,\%\\
\bottomrule
\end{tabular}}
\par\vspace{4pt}
{\small The \textbf{leader} is simply the candidate with the best \emph{observed} validation score:
here \dat{c1}. It is not automatically the final answer.}
\end{minipage}

% ---------------------------------------------------------------- S5
\step{S5}{Paired bootstrap --- is a difference real or noise?}\tagval

\begin{center}
\begin{tikzpicture}[node distance=6mm]
  \node[valid, text width=3.2cm, font=\footnotesize] (p1) {same validation row $i$\\$d_i = \text{score}_A(i) - \text{score}_B(i)$};
  \node[code, text width=2.9cm, font=\footnotesize, right=of p1] (p2) {$d_1, d_2, \dots, d_{1575}$\\observed difference\\$\hat\Delta = \text{mean}(d)$};
  \node[code, text width=2.9cm, font=\footnotesize, right=of p2] (p3) {resample the rows\\with replacement,\\\textbf{2\,000 times}};
  \node[code, text width=2.6cm, font=\footnotesize, right=of p3] (p4) {middle 95\,\% of the\\2\,000 means\\= \textbf{95\,\% CI}};
  \draw[arr] (p1) -- (p2); \draw[arr] (p2) -- (p3); \draw[arr] (p3) -- (p4);
\end{tikzpicture}
\end{center}

{\small Pairing by row matters: rows that are easy or hard for \emph{every} model cancel out, so the interval
reflects how the two candidates actually differ. The uncertainty comes from the finite set of validation
rows --- not from the seeds. All of this is numpy code; no LLM is involved.}

\vspace{6pt}
\textbf{Three comparisons, three different jobs}\par\vspace{3pt}
{\small
\begin{tabularx}{\linewidth}{@{}l Y Y@{}}
\toprule
\textbf{comparison} & \textbf{question it answers} & \textbf{what it is used for}\\ \midrule
\textbf{vs Leader} & Is this candidate clearly worse than the best one observed so far? & \textbf{decides contender status} --- the only one that does\\
vs Reference & How does it compare with the original baseline? & context for the Recommender and the UI\\
vs Parent & Did the latest one-knob change (parent $\rightarrow$ child) help? & evidence for the next refinement (refined candidates only)\\
\bottomrule
\end{tabularx}}

\vspace{8pt}
\begin{minipage}[t]{0.47\linewidth}
\vspace{0pt}
\textbf{The contender rule (vs Leader)}\par\vspace{4pt}
\begin{tikzpicture}[node distance=3mm]
  \node[valid, text width=3.4cm, font=\footnotesize] (c) {CI of (candidate $-$ leader)};
  \node[code, text width=3.0cm, font=\footnotesize, below left=6mm and -1.2cm of c] (w) {entirely below 0\\$\rightarrow$ clearly worse\\$\rightarrow$ \textbf{eliminated}};
  \node[code, text width=3.0cm, font=\footnotesize, below right=6mm and -1.2cm of c, draw=val, line width=1pt] (k) {overlaps 0\\$\rightarrow$ not clearly worse\\$\rightarrow$ \textbf{contender}};
  \draw[arr] (c) -- (w); \draw[arr] (c) -- (k);
\end{tikzpicture}
\par\vspace{3pt}
\note{The leader is always a contender. For MSE (lower is better) ``worse'' means entirely above 0.
An eliminated family is never refined.}
\end{minipage}\hfill
\begin{minipage}[t]{0.49\linewidth}
\vspace{0pt}
\textbf{Real run, after round 1} \note{(vs leader c1, percentage points)}\par\vspace{3pt}
{\small
\begin{tabular}{@{}l r l@{}}
\toprule
candidate & difference & 95\,\% CI $\rightarrow$ status\\ \midrule
c1 linear & --- & leader $\rightarrow$ \textbf{contender}\\
c2 tree & $-3.30$ & $[-5.08, -1.59]$ $\rightarrow$ eliminated\\
c3 forest & $-0.85$ & $[-1.90, +0.25]$ $\rightarrow$ \textbf{contender}\\
c4 MLP & $-1.65$ & $[-2.88, -0.38]$ $\rightarrow$ eliminated\\
\bottomrule
\end{tabular}}
\par\vspace{4pt}
\note{The forest is only 0.85 points behind and that gap is within noise, so it stays in the race.
The MLP is only 1.65 points behind, but its interval never reaches 0.}
\end{minipage}

% ---------------------------------------------------------------- S6
\step{S6}{Decide --- stop, or ask the Recommender?}{\tagcode\ \tagrec}

\begin{minipage}[t]{0.47\linewidth}
\vspace{0pt}
\begin{tikzpicture}[node distance=5mm]
  \node[valid, text width=3.6cm, font=\footnotesize] (v) {validation analysis\\(leader, contenders)};
  \node[code, text width=3.6cm, font=\footnotesize, below=of v] (m) {Is this round 3\\(\dat{MAX\_ROUNDS})?};
  \node[code, text width=3.6cm, font=\footnotesize, below=of m] (f) {Are all contenders the\\same model family?};
  \node[recom, text width=3.6cm, font=\footnotesize, below=of f] (r) {ask the \textbf{Recommender}};
  \node[code, text width=1.7cm, font=\footnotesize, right=7mm of m] (c1) {CONCLUDE\\{\scriptsize budget}};
  \node[code, text width=1.7cm, font=\footnotesize, right=7mm of f] (c2) {CONCLUDE\\{\scriptsize settled}};
  \node[code, text width=1.7cm, font=\footnotesize, right=7mm of r] (c3) {REFINE or\\CONCLUDE};
  \draw[arr] (v) -- (m);
  \draw[arr] (m) -- node[left, lab] {no} (f);
  \draw[arr] (f) -- node[left, lab] {no} (r);
  \draw[arr] (m) -- node[above, lab] {yes} (c1);
  \draw[arr] (f) -- node[above, lab] {yes} (c2);
  \draw[arr] (r) -- (c3);
\end{tikzpicture}
\end{minipage}\hfill
\begin{minipage}[t]{0.49\linewidth}
\vspace{0pt}
{\small The two stop rules are code, checked before any LLM call. ``Settled'' means tuning can no longer
change \emph{which family} is best. In the real run the contenders after round 1 were linear and forest
--- two families --- so the Recommender was asked.}
\par\vspace{5pt}
{\small
\begin{tabularx}{\linewidth}{@{}Y Y@{}}
\toprule
\textbf{The Recommender sees} & \textbf{It never sees}\\ \midrule
each candidate's family, parent and change; validation and train score; overfitting gap; seed spread;
contender flag; differences + CIs vs reference, leader and parent; the allowed knob ranges of the
contenders' families; rounds left
& raw dataset rows; per-row score arrays; anything from the test set; earlier conversations\\
\bottomrule
\end{tabularx}}
\par\vspace{3pt}
\note{It receives these as one compact JSON object --- about 730 tokens in the real run.}
\end{minipage}

% ---------------------------------------------------------------- S7
\needspace{20\baselineskip}
\step{S7}{Refine --- one contender, one knob, 1--3 values}{\tagrec\ \tagcode}

\begin{minipage}[t]{0.5\linewidth}
\vspace{0pt}
\textbf{What the Recommender may change}
\begin{enumerate}
\item pick \textbf{one} parent --- it must be a contender;
\item pick \textbf{one} knob of that parent's family (e.g.\ forest: \dat{max\_depth},
      \dat{min\_samples\_leaf}, \dat{max\_features});
\item propose \textbf{1--3 values}, each inside the knob's allowed range.
\end{enumerate}
\par\vspace{4pt}
\textbf{What code checks before anything runs}
\begin{itemize}
\item all three rules above;
\item no new configuration was already tried;
\item a violation is sent back to the LLM once; an invalid second answer fails the investigation.
\end{itemize}
\par\vspace{4pt}
\note{Linear regression has no tunable knobs, so after round 1 the only refinable contender was the forest.
A request to refine the MLP (c4) is rejected: ``c4 is not a contender''.}
\end{minipage}\hfill
\begin{minipage}[t]{0.46\linewidth}
\vspace{0pt}
\textbf{Example} \note{(illustration with the maximum of 3 values)}\par\vspace{2pt}
{\small parent = c3 forest, knob = \dat{min\_samples\_leaf}, values = 2, 4, 8}\par\vspace{4pt}
\begin{tikzpicture}[x=1cm, y=-1cm]
  \node[code, font=\footnotesize] (p) at (2.1,0) {c3 forest (parent)};
  \node[code, font=\footnotesize] (k1) at (0.55,1.15) {leaf=2};
  \node[code, font=\footnotesize] (k2) at (2.1,1.15) {leaf=4};
  \node[code, font=\footnotesize] (k3) at (3.65,1.15) {leaf=8};
  \foreach \k in {k1,k2,k3} \draw[arr] (p) -- (\k);
  \foreach \k/\x in {k1/0.55,k2/2.1,k3/3.65} {
    \node[run] (\k a) at (\x,1.95) {seed 0};
    \node[run] (\k b) at (\x,2.45) {seed 1};
    \node[run] (\k c) at (\x,2.95) {seed 2};
    \draw[arr] (\k) -- (\k a);
  }
  \node[lab] at (5.0,1.15) {3 new};
  \node[lab] at (5.0,1.45) {candidates};
  \node[lab] at (5.0,2.45) {9 runs};
\end{tikzpicture}
\par\vspace{4pt}
{\small Each child is the parent with \emph{only} \dat{min\_samples\_leaf} changed.
\textbf{3 candidates $\times$ 3 seeds = 9 runs}, then back to S3: train $\rightarrow$ validation
$\rightarrow$ bootstrap $\rightarrow$ new contenders $\rightarrow$ decide again.}
\end{minipage}

\vspace{8pt}
\textbf{What the real run did} \note{(each child's ``vs parent'' difference is the effect of one knob)}\par\vspace{4pt}
\begin{center}
\begin{tikzpicture}[x=1cm, y=-1cm]
  \node[code, font=\footnotesize] (c3) at (0,0) {c3 forest};
  \node[code, font=\footnotesize] (c5) at (3.6,-0.35) {c5\, leaf=5};
  \node[code, font=\footnotesize] (c6) at (3.6,0.45) {c6\, leaf=20};
  \node[code, font=\footnotesize, draw=val, line width=1pt] (c7) at (7.6,-0.75) {c7\, + max\_depth=12};
  \node[code, font=\footnotesize] (c8) at (7.6,0.05) {c8\, + max\_depth=24};
  \draw[arr] (c3) -- (c5); \draw[arr] (c3) -- (c6);
  \draw[arr] (c5) -- (c7); \draw[arr] (c5) -- (c8);
  \node[lab] at (1.8,-1.55) {after round 1: refine c3};
  \node[lab] at (5.6,-1.55) {after round 2: refine c5};
  \node[lab, align=left, anchor=west] at (9.5,-0.35) {after round 3:\\stop (budget)};
\end{tikzpicture}
\end{center}
{\small Round 2: 2 values $\times$ 3 seeds = 6 runs. Round 3: 6 runs. Total: 8 candidates, 20 runs. All
forest children stayed contenders; none was clearly better than its parent (e.g.\ c7 vs c5:
$+0.08$ points, CI $[-0.38, +0.53]$).}

% ---------------------------------------------------------------- S8
\step{S8}{Finalize --- the only time the test set is opened}{\tagcode\ \tagtst}

\begin{minipage}[t]{0.52\linewidth}
\vspace{0pt}
\begin{enumerate}
\item \textbf{Winner $W$} (validation only): the best candidate \emph{other than the reference}.
      Normally that is the final leader; if the reference itself leads, $W$ is its strongest alternative.
\item \textbf{Runner-up $R$} (validation only): the best candidate whose family is neither $W$'s nor the
      reference's. If none exists, there is no second comparison.
\item \textbf{Only now} the sealed test scores of $W$, the reference and $R$ are read.
\item The same paired bootstrap, on the 1\,575 test rows:
      \textbf{$W$ vs reference} and \textbf{$W$ vs $R$}.
\end{enumerate}
\par\vspace{5pt}
With $k$ final comparisons, each CI uses (Bonferroni)
\[
\alpha = \frac{0.05}{k}:\qquad k = 2 \;\Rightarrow\; \alpha = 0.025,\; \text{confidence } 97.5\,\%
\]
\note{$k = 1$ (no runner-up) $\Rightarrow$ 95\,\%. The validation intervals in S5 are not corrected:
they only choose experiments, they are never reported as the answer.}
\end{minipage}\hfill
\begin{minipage}[t]{0.44\linewidth}
\vspace{0pt}
\begin{tikzpicture}
  \node[test, text width=4.6cm, font=\footnotesize] (t) {\textbf{TEST SET}};
  \node[lab, below=1mm of t, text width=4.6cm, align=center] (n) {never used during the adaptive loop};
  \node[code, text width=4.6cm, font=\footnotesize, below=5mm of n] (f) {finalize only, comparisons fixed first};
  \draw[arr, tst, dashed] (n) -- (f);
\end{tikzpicture}
\par\vspace{6pt}
\textbf{Real run} \note{(97.5\,\% CIs, percentage points)}\par\vspace{3pt}
{\small
\begin{tabular}{@{}l l@{}}
\toprule
$W$ & c7 forest (79.6\,\% on test)\\
reference & c1 linear (80.8\,\% on test)\\
$R$ & c4 MLP\\ \midrule
$W$ vs reference & $-1.2$ $[-2.3, -0.1]$ \textbf{worse}\\
$W$ vs $R$ & $+0.2$ $[-1.1, +1.6]$ \textbf{tied}\\
\bottomrule
\end{tabular}}
\par\vspace{4pt}
\note{Answer: plain logistic regression beat the tuned forest; the forest and the MLP are statistically
tied. The report also lists 4 LLM calls, 20 runs, and how many candidates each family got.}
\end{minipage}

% ============================================================================
%  PART 4 -- EFFECT MODE
% ============================================================================
\newpage
\partx{4}{Effect Mode --- ``does factor X matter?''}

{\small Running example (real run): \emph{``Does tree depth matter for a decision tree on this dataset?''}
The Planner returns \dat{factor = max\_depth}, \dat{levels = 2, 4, 8, 16}, \dat{reference = 8},
base model: decision tree.}

\vspace{4pt}
\begin{keybox}
Effect Mode does \textbf{not} compare model families. Everything is held fixed except one factor, and the
question is: \emph{does changing this factor away from the reference level make a measurable difference?}
The factor can be \dat{model\_type}, \dat{normalize}, or any hyperparameter of the base family.
\end{keybox}

\begin{center}
\begin{tikzpicture}[node distance=4.2mm]
  \node[code, minimum width=6.4cm] (e1) {E1\; one candidate per level};
  \node[code, minimum width=6.4cm, below=of e1] (e2) {seeds $\rightarrow$ runs, train (as S2--S3)};
  \node[valid, minimum width=6.4cm, below=of e2] (e3) {E2\; each level vs the reference level};
  \node[code, minimum width=6.4cm, below=of e3] (e4) {stop rule: round 3 reached? (code)};
  \node[code, minimum width=6.4cm, below=9mm of e4] (e5) {E4\; choose the best non-reference level};
  \node[test, minimum width=6.4cm, below=of e5] (e6) {E4\; open test: level vs reference ($k = 1$)};
  \foreach \a/\b in {e1/e2,e2/e3,e3/e4,e5/e6} \draw[arr] (\a) -- (\b);
  \draw[arr] (e4) -- node[right, lab] {stop} (e5);
  \node[recom, text width=3.6cm, right=12mm of e4] (r) {E3\; \textbf{Recommender}:\\{\footnotesize 1--3 new levels of\\the same factor}};
  \node[code, text width=3.6cm, right=12mm of e2] (b) {code validates, builds\\new levels from the base};
  \draw[arr] (e4) -- node[above, lab] {continue} (r);
  \draw[arr] (r) -- (b);
  \draw[arr] (b) -- node[above, lab] {new} (e2);
  \draw[arr] (r.south) |- node[pos=0.25, right, lab] {conclude} (e5.east);
\end{tikzpicture}
\end{center}

\step{E1}{Initial candidates --- one per level}\tagcode

\begin{minipage}[t]{0.45\linewidth}
\vspace{0pt}
\begin{tikzpicture}[x=1cm, y=-1cm]
  \node[code, font=\footnotesize, draw=plan, line width=1pt] (r) at (0,1.2) {c1\; max\_depth=8\\{\scriptsize reference}};
  \node[code, font=\footnotesize] (a) at (3.6,0)   {c2\; max\_depth=2};
  \node[code, font=\footnotesize] (b) at (3.6,0.8) {c3\; max\_depth=4};
  \node[code, font=\footnotesize] (c) at (3.6,1.6) {c4\; max\_depth=16};
  \foreach \x in {a,b,c} \draw[arr] (r.east) -- (\x.west);
  \node[lab] at (1.8,2.35) {only max\_depth differs};
\end{tikzpicture}
\end{minipage}\hfill
\begin{minipage}[t]{0.51\linewidth}
\vspace{0pt}
\begin{panel}
\small\textbf{c3 compared with c1}\par\vspace{2pt}
same model (decision tree) $\cdot$ same preprocessing $\cdot$ same other hyperparameters\\
\textbf{except} \dat{max\_depth = 4} instead of 8
\end{panel}
{\small That is what makes the comparison interpretable: any difference is caused by the factor.
A decision tree is seed-independent, so each level is \textbf{1 run} (4 runs). An MLP factor would
give 3 runs per level.}
\end{minipage}

\step{E2}{Run and analyse --- the same machinery, a simpler comparison}\tagval

{\small Training, row scores, the sealed test set, seed averaging and the paired bootstrap are exactly
S2--S5. The difference: every level is compared with the \textbf{reference level} only. There are no
leader comparisons, no contenders and no ``settled'' rule in Effect Mode.}

\vspace{5pt}
\begin{minipage}[t]{0.49\linewidth}
\vspace{0pt}
{\small
\begin{tabular}{@{}l r l@{}}
\toprule
\textbf{level vs 8} & \textbf{val} & \textbf{difference, 95\,\% CI}\\ \midrule
8 (reference) & 75.8\,\% & ---\\
2 & 73.9\,\% & $-1.8$ $[-3.8, +0.2]$\\
4 & 75.6\,\% & $-0.2$ $[-2.0, +1.5]$\\
16 & 71.6\,\% & $-4.1$ $[-6.0, -2.2]$ \textbf{worse}\\
\bottomrule
\end{tabular}}
\par\vspace{2pt}
\note{Real run, after round 1, percentage points.}
\end{minipage}\hfill
\begin{minipage}[t]{0.47\linewidth}
\vspace{0pt}
{\small The question for each row: \emph{does moving the factor from 8 to this level produce a
measurable difference?} Depth 16 clearly hurts (it scores 98\,\% on its training rows: overfitting);
2 and 4 are within noise of 8.}
\end{minipage}

\step{E3}{Refine --- more values of the same factor}{\tagrec\ \tagcode}

\begin{minipage}[t]{0.5\linewidth}
\vspace{0pt}
{\small The Recommender sees the same kind of compact summary and may only propose \textbf{1--3 new values
of the same factor}. Code forces the knob to be the factor and the parent to be the reference: every new
level is built from the same base configuration, never from another level.}
\par\vspace{4pt}
\begin{itemize}
\item[] \small after round 1: \emph{``try 6 and 12''} $\rightarrow$ c5, c6
\item[] \small after round 2: \emph{``try 5 and 7''} $\rightarrow$ c7, c8
\item[] \small after round 3: stop (budget)
\end{itemize}
\end{minipage}\hfill
\begin{minipage}[t]{0.46\linewidth}
\vspace{0pt}
\begin{tikzpicture}[x=1cm, y=-0.6cm]
  \node[code, font=\footnotesize, draw=plan, line width=1pt] (r) at (0,1.5) {c1 depth 8\\{\scriptsize reference}};
  \foreach \l/\y/\rd in {2/-0.5/1, 4/0.25/1, 16/1.0/1, 6/1.75/2, 12/2.5/2, 5/3.25/3, 7/4.0/3}
    { \node[run, anchor=west] (n\l) at (2.0,\y) {depth \l}; \draw[arr] (r.east) -- (n\l.west);
      \node[lab, anchor=west] at (3.3,\y) {round \rd}; }
\end{tikzpicture}
\par\vspace{2pt}
\note{Every level hangs off the reference: one factor, one star.}
\end{minipage}

\step{E4}{Finalize --- one confirmatory comparison}{\tagcode\ \tagtst}

\begin{minipage}[t]{0.5\linewidth}
\vspace{0pt}
\begin{enumerate}
\item From validation: the best non-reference level --- here depth 6 (77.0\,\%).
\item Open the sealed test set once.
\item One comparison: \textbf{selected level vs reference level}, so $k = 1$, $\alpha = 0.05$,
      \textbf{95\,\%} --- no Bonferroni division is needed.
\end{enumerate}
\end{minipage}\hfill
\begin{minipage}[t]{0.46\linewidth}
\vspace{0pt}
\begin{panel}
\small\textbf{Real run --- test, 1\,575 rows}\par\vspace{2pt}
depth 6: 78.9\,\% $\cdot$ depth 8: 77.8\,\%\\
difference $+1.1$ points, 95\,\% CI $[-0.3, +2.4]$\\
$\rightarrow$ \textbf{no reliable difference}\\[2pt]
8 runs in 4 s; stopped by the round budget.
\end{panel}
\end{minipage}

% ============================================================================
%  PART 5 -- SIDE BY SIDE
% ============================================================================
\vspace{6pt}
\partx{5}{The two modes side by side}

{\small
\begin{tabularx}{\linewidth}{@{}l Y Y@{}}
\toprule
 & \textbf{Selection Mode} & \textbf{Effect Mode}\\ \midrule
Question & Which model family is best? & Does factor X matter?\\
Initial design & 2--4 model families, each at its defaults & 1 factor, 2--5 levels, a base configuration\\
Reference & the simplest chosen family (code picks it) & the reference level the Planner names\\
Seeds & 3 runs per forest/MLP candidate, 1 per linear/tree candidate & same\\
Validation comparisons & vs reference, vs leader, vs parent & vs reference only\\
Contenders & yes: not clearly worse than the leader & none\\
Stop rules (code) & round 3 reached, or all contenders one family & round 3 reached\\
Refinement & one contender, one knob of its family, 1--3 values & 1--3 new levels of the same factor\\
Parent of a new candidate & the contender it refines (can be a refined child) & always the reference level\\
Final comparisons & $W$ vs reference, $W$ vs runner-up from another family & best level vs reference\\
Test confidence & 97.5\,\% each if 2 comparisons (95\,\% if 1) & 95\,\%\\
\bottomrule
\end{tabularx}}

\vspace{8pt}
{\small\textbf{Hard limits (both modes, all enforced by code):} at most 3 rounds $\cdot$ at most 3 new
candidates per refinement $\cdot$ MLP epochs $\le 50$, batch size $\ge 16$ $\cdot$ at most 4 LLM calls
(Planner, 2 Recommender decisions, 1 explanation) $\cdot$ at most 26 runs (Selection: $8 + 2 \times 3
\times 3$) or 33 runs (Effect: $(5 + 2 \times 3) \times 3$).
\textbf{Limitation:} every result is conditional on the one fixed training split.}

\end{document}
